\documentclass[12pt,letterpaper]{article}

% --- Configuración de Idioma y Tipografía Industrial (Celtx) ---
\usepackage[T1]{fontenc} 
\usepackage[utf8]{inputenc}
\usepackage[spanish]{babel}
\usepackage{courier} 
\renewcommand{\familydefault}{\ttdefault} % Todo en Courier

% --- Alineación Izquierda de Guion ---
\usepackage{ragged2e} 

% --- Márgenes Estándar de la Industria ---
\usepackage{geometry}
\geometry{
    letterpaper,
    left=1.5in,   
    right=1.0in,
    top=1.0in,
    bottom=1.0in
}

% --- Espaciado ---
\usepackage{setspace}
\setstretch{1.15} 
\parindent 0in    

% --- Contador Automático de Escenas ---
\newcounter{numescena}
\setcounter{numescena}{0}

% --- Definición de Comandos de Entrada (Variables) ---
\newcommand{\duration}[1]{\gdef\valduration{#1}}
\newcommand{\scenetype}[1]{\gdef\valscenetype{#1}}
\newcommand{\actors}[1]{\gdef\valactors{#1}}
\newcommand{\notes}[1]{\gdef\valnotes{#1}}
\newcommand{\desc}[1]{\gdef\valdesc{#1}}

% --- Inicializador por defecto (Para evitar errores si se deja alguno vacío) ---
\newcommand{\limpiardatos}{%
    \gdef\valduration{---}%
    \gdef\valscenetype{---}%
    \gdef\valactors{---}%
    \gdef\valnotes{---}%
    \gdef\valdesc{}%
}

% --- Entorno Principal de la Escena ---
% Uso: \begin{scene}{INT/EXT. LUGAR}{MOMENTO} ... \end{scene}
\newenvironment{scene}[2]{%
    \refstepcounter{numescena}%
    \limpiardatos % Resetea las variables anteriores
    \gdef\valhead{#1 - #2}% Guarda el encabezado
}{%
    % Aquí se ejecuta la impresión real respetando el orden del template CEDEC
    \vspace{2.0em} 
    \noindent\textbf{\arabic{numescena}. \valhead}%
    \nopagebreak\vspace{0.4em}\par
    
    \noindent\textbf{DURACIÓN:} \valduration\par
    \nopagebreak\noindent\textbf{TIPO DE ESCENA:} \valscenetype\par
    \nopagebreak\noindent\textbf{PERSONAJES:} \valactors\par
    \nopagebreak\noindent\textbf{OTROS DATOS:} \valnotes\par
    \nopagebreak\vspace{0.6em}\par 
    
    \noindent\textbf{DESCRIPCIÓN:}\\\relax
    \valdesc\par
}

% --- PORTADA ESTILO CELTX ---
\begin{document}
\RaggedRight

\thispagestyle{empty}
\vspace*{2.5in}
\begin{center}
    {\Huge \textbf{[SIN TÍTULO]}}\\[1cm]
    {\large DRAFT: ESCALETA}\\[2cm]
    Escrito por\\
    \textbf{[CINCOSEIS]}
\end{center}

\vfill
\noindent SERGIO GONZÁLEZ\\
Contacto: [ Email / Teléfono]\hfill \today
\newpage

% --- INICIO DE LA ESCALETA ---
\pagestyle{plain}
\setcounter{page}{1}

\textbf{SINOPSIS} \\
KSJDKSDJKSD
JASHDJAHDJKASD



ASHJDGJKASGDHJKAGHJKDGHAJSDGHJSA
\section*{\centering ACTO I}


% EJEMPLO 1: Escrito en el orden sugerido
\begin{scene}{INT. VESTÍBULO DEL TEATRO}{DÍA}
    \duration{1' 30''}
    \scenetype{HABLADA / CANTADA}
    \actors{Carlos, Directora}
    \notes{Efecto de eco de teatro vacío, partituras sobre el piano.}
    \desc{Carlos (25) se presenta a las audiciones de la nueva obra musical. Tiembla de los nervios. La Directora, sentada en la penumbra del patio de butacas, le da una señal seca con la mano.}
\end{scene}

% EJEMPLO 2: El orden no importa. Escribimos primero la descripción y al final los actores.
\begin{scene}{INT. ESCENARIO - MÁS TARDE}{DÍA}
    \desc{Despechado por su fallo anterior, Carlos pide una segunda oportunidad para demostrar sus habilidades de baile. El ensamble musical arranca con un ritmo de zarzuela acelerado.}
    \notes{Cambio de iluminación a tonos azules, piano en directo.}
    \scenetype{COREOGRAFIADA / TOCADA}
    \duration{2' 15''}
    \actors{Carlos, Elenco de baile, Músicos}
\end{scene}

\end{document}